\section{Startup Sequence}
\label{sec:rec_init}
\label{sec:startup}

When the \texttt{micro:bit v1} device starts, it reads the initial stack pointer from
address \lstinline[style=myagc]|0x0| into the \lstinline[style=myagc]|sp| register. It
then begins executing the code at the instruction address stored at
\lstinline[style=myagc]|0x4|. The \lstinline[style=myagc]|__stack| and
\lstinline[style=myagc]|__reset| symbols are placed at the start of the vector table
along with the interrupt handlers. This table is loaded at address
\lstinline[style=myagc]|0x0| by the linker script.

\begin{lstlisting}[
    style=myagc,
    caption={\file{src/vector\_table.c}},
]
__attribute((section(".vectors"))) void *__vectors[] = {
    __stack, // Initial Main Stack Pointer (MSP) value
    __reset, // Reset Handler: Entry point after power-up or reset

    ... // interrupt handlers
};
\end{lstlisting}

It is the job of the \lstinline[style=myagc]|__reset| function to initialize the clock,
set up memory sections, initialize hardware and drivers, run tests, and finally start
the scheduler.

\begin{lstlisting}[
    style=myagcs,
    caption={\file{src/startup.cpp}},
    label={code:startup},
    emph={[6] __extern_C__ },
    emphstyle={[6]\typestyle},
    emph={[7] SEC_ZERO, is_boot, clear_screen, trigger_hardfault, call_constructors},
    emphstyle={[7]\stdlibstyle},
    linebackgroundcolor={\lines{6}{10} \line{12} \lines{14}{17} \line{19} \lines{21}{23}},
]
__extern_C__ void __reset(void)
{
  /* Activate the crystal clock */
  /* Set up write-back protection for the first half of FLASH */

  SEC_INIT(data);         
  SEC_ZERO(bss);          +-------------------------------+
  SEC_INIT(services);     | Initialize data sections      |
  SEC_INIT(startups);     +-------------------------------+
  SEC_INIT(apps);

  call_constructors();    // Manually call necessary global constructors

  led::init();            +-------------------------------+
  serial::init();         | Set up hardware and drivers   |
  i2c::init();            +-------------------------------+
  fs::init();             

  boot::init();           // Identify the type of boot

  recover::init();        +-------------------------------+
  tests::run();           | Run tests and start processes |
  sched::init();          +-------------------------------+

  // Should never get here
  trigger_hardfault();
}
\end{lstlisting}

\paragraph*{Initialization of Data Sections} 

During linking, the linker places all static and global data structures into specific
sections in FLASH. These must be copied to RAM at startup at pre-determined addresses.
Thus, the first task of \lstinline[style=myagc]|__reset| is to initialize these
sections. The macro \lstinline[style=myagc]|SEC_INIT| copies data from FLASH to RAM,
while \lstinline[style=myagc]|SEC_ZERO| zero-initializes it.

Besides the default \lstinline[style=myagc]|.bss| and \lstinline[style=myagc]|.data|
sections, we define additional custom sections:
\begin{itemize}
    \item \lstinline[style=myagc]|.recover|: Objects in this section are retained
        across soft resets.
    \item \lstinline[style=myagc]|.services|: Stores definitions for service processes
        (e.g., \lstinline[style=myagc]|display|, see \autoref{proc:display}) which are
        started at every boot.
    \item \lstinline[style=myagc]|.startups|: Contains processes only launched after a
        hardware reset, but not during recoveries.
    \item \lstinline[style=myagc]|.apps|: Contains processes launchable via the shell
        (\autoref{sec:shell}).
\end{itemize}

The testing framework also defines \lstinline[style=myagc]|.unit_tests| and
\lstinline[style=myagc]|.sys_tests| sections for storing test objects
(\autoref{sec:testing}).

\paragraph*{Dynamic Global Object Construction} 

In embedded settings, global object constructors are usually disabled for performance.
However, some non-trivial constructors must be run at startup. To support this, we
manually define initializers in a special FLASH section
\lstinline[style=myagc]|.init_funcs|, and call them explicitly via
\lstinline[style=myagc,
basicstyle=\bfseries\codestyle\codefontsize]|call_constructors|.

\paragraph*{Driver Initialization} 

\lstinline[style=myagc]|__reset| then initializes modules and device drivers. These
\lstinline[style=myagc]|init| functions set up hardware peripherals and internal data
structures as needed.

\paragraph*{Boot Type Identification} 

As discussed in \autoref{sec:rec_tbl}, recovery mechanisms depend on knowing the boot
type. We distinguish among:
\begin{itemize}
    \item \texttt{FLASH} — Full flash programming occurred.
    \item \texttt{POWER} — Power was disconnected and reconnected.
    \item \texttt{BOOT} — Reset button triggered restart.
    \item \texttt{RESET} — Software-triggered soft reset.
\end{itemize}

Detection relies on:
\begin{itemize}
    \item \texttt{FLASH} — Missing magic value in flash indicates reprogramming.
    \item \texttt{RESET} — RAM-retained magic value indicates a soft reset.
    \item \texttt{POWER} — Zeroed \lstinline[style=myagc]|POWER.RESETREAS| register
        indicates power-on reset.
    \item \texttt{BOOT} — Default fallback if none of the above apply.
\end{itemize}

\paragraph*{Boot Statistics} 

We track and store boot type statistics in a FRAM file, printed at every startup for
diagnostic (and amusement) purposes:
\begin{lstlisting}[style=mytxt]
boot level: flash
boot stat: n_flash = 395, n_boot = 22, n_power = 19, n_reset = 114
\end{lstlisting}

This output shows the current boot level (\texttt{flash}) and counts for different reset
modes since the last FRAM filesystem format.


\subsection{Scheduler Initialization}

After memory, hardware devices, and drivers have been initialized, the recovery table
is set up:
\begin{itemize}
    \item If the boot type is \lstinline[style=myagc]|FLASH| or
        \lstinline[style=myagc]|BOOT|, the recovery table is cleanly initialized.
    \item Otherwise, the table is copied from its file stored in FRAM into RAM.
    \item If the boot type is not \lstinline[style=myagc]|RESET|, the
        \lstinline[style=myagc]|.recover| section is also copied from FLASH to RAM.
\end{itemize}

If there are tests to run, they are executed at this stage (see \autoref{sec:testing}
for details).

Finally, the scheduler is initialized:
\begin{itemize}
    \item The idle process is registered (\autoref{sec:idle_proc}).
    \item Service processes are registered.
    \item Recovery routines are executed, and startup processes are initialized.
    \item The system timer is started.
    \item Control transfers to the idle process, which immediately yields to the
        highest-priority process.
\end{itemize}
This completes system startup.

\begin{lstlisting}[
    style=myagc,
    caption={\file{src/core/sched.cpp}},
    linebackgroundcolor={\lines{6}{7} \line{10} \line{13} \line{17} \line{21}
    \lines{29}{30}},
    emph={[7]setup, PROC_FUNC, intr_disable, IDLE_PID, setup_startups, recover},
    emphstyle={[7]\stdlibstyle},
    label={sched:init},
    mathescape=true,
]
/* assign idle_task to the main process, 
 * sets up service processes, 
 * then enters the idle_task */
void init()
{
  intr_disable(); 
    // disable all interrupts, which will be enabled by the idle proc later

  // register the idle process, and mark it as the currently running process
  IDLE_PID = REG_PROC(idle, nullptr);
  cpu.curr_proc = procs[IDLE_PID];

  setup_services(); // register all service processes

  if (/* if boot type is FLASH or BOOT */) {
    // setup the startup processes from .startup section
    sched::setup_startups();
  } else {
    // run the recovery mechanisms set in the recovery table
    auto reset_lev = /* RESET or POWER_OFF */;
    recover(reset_lev); // $\text{\autoref{code:recover}}$
  }

  // start the timer
  timer::init();

  // then jump to the execution of the idle task by a context switch
  // which as seen in $\text{\autoref{sec:idle_proc}}$, immediately calls change_proc
  cpu.curr_proc->stk_ptr = cpu.curr_proc->stack + cpu.curr_proc->stk_sz;
  __run(PROC_FUNC(idle), &cpu.curr_proc->stk_ptr);
}
\end{lstlisting}

